\chapter{Project Context and Problem Definition}

\section{Introduction}

Cybersecurity monitoring usually concentrates on production resources: user
accounts, servers, business applications, and sensitive data. This produces a
large volume of legitimate activity that analysts must distinguish from
malicious behavior. Cyber deception changes this asymmetry by introducing
resources that have no normal business use. When such a resource is accessed,
the event deserves immediate investigation. NIST includes deception and
honeytokens among the techniques that can support cyber-resilient systems
\cite{nist800160v2}.

JANUS focuses on a specific deception resource: the honey document. Examples
include a financial workbook, an internal procedure, a technical configuration,
or decoy cloud credentials. The file must be sufficiently plausible to attract
an intruder, while remaining harmless and distinguishable from real data by the
defender.

\section{Problem statement}

A static decoy file is easy to deploy, but it presents four practical
limitations.

\begin{enumerate}[label=\arabic*.]
  \item \textbf{Limited realism and diversity.} Reusing the same names and
  content makes a deception environment easier to fingerprint.
  \item \textbf{Incomplete observation.} A remote callback can show that a
  token endpoint was contacted, but it does not by itself explain which local
  account and process touched the file.
  \item \textbf{Format-dependent activation.} A Word beacon, an Excel token,
  a web breadcrumb, and AWS decoy credentials do not activate in the same way.
  \item \textbf{Risk of unsupported conclusions.} An event may identify a
  process or an egress IP without proving a person's identity, a device MAC
  address, or a complete attack path.
\end{enumerate}

The central question is therefore:

\begin{quote}
How can a defensive platform generate credible and controlled honey documents,
deploy them safely, and combine local and remote evidence to detect attacker
interaction without fabricating unavailable attributes?
\end{quote}

\section{Research questions}

The project is structured around the following questions:

\begin{itemize}
  \item How can an LLM vary business content without receiving local paths,
  secrets, or live token material?
  \item Which Canarytoken type is appropriate for each generated file format?
  \item Which Windows and Wazuh fields are direct evidence of file access?
  \item How can raw Wazuh evidence, normalized events, and verdicts remain
  traceable to one HoneyDoc instance?
  \item Which claims about the attacker are reliable, inferred, or impossible?
\end{itemize}

\section{Objectives}

\subsection{Primary objective}

The primary objective is to implement a reproducible MVP that joins four
capabilities: AI-assisted content generation, format-specific tokenization,
controlled deployment, and evidence correlation.

\subsection{Functional objectives}

\begin{itemize}
  \item Generate financial, human-resources, technical, and cloud-credential
  decoys in DOCX, XLSX, CSV, JSON, YAML, ENV, and ZIP formats.
  \item Create official Canarytokens only after the content passes safety and
  coherence checks.
  \item Deploy files atomically under a single authorized root and record their
  SHA-256 hashes in SQLite.
  \item Detect local read, modification, and deletion events using Windows
  auditing and Wazuh rules.
  \item Preserve raw Wazuh alerts and expose normalized fields through an
  authenticated API and dashboard.
  \item State token activation conditions and sensor limitations in the user
  interface.
\end{itemize}

\subsection{Security and quality objectives}

The system must avoid macros by default, prevent path traversal, keep
administrative routes behind an API key, avoid printing secret token material,
and fail closed when the configured Canarytokens provider is unavailable.
Tests must cover artifact structure, database migration, API security,
deployment policy, token-provider behavior, Wazuh parsing, verdict generation,
and dashboard navigation.

\section{Scope and non-goals}

The supported core contains the HoneyDoc pipeline, Wazuh, Canarytokens, SQLite,
FastAPI, and Streamlit. Experimental fake SSH and HTTP services are excluded
from the validated architecture because they depend on a later attacker choice
to reuse exposed credentials; they do not improve the initial observation of a
document interaction. Prompt-injection and credential-trap layers also remain
disabled by default.

The project is not a production SOC platform. It does not provide high
availability, immutable storage, global endpoint coverage, attribution of a
human identity, or retrieval of a remote MAC address. Wazuh observes only hosts
on which an agent and the required audit policy are active. A Canarytokens hit
can continue to work after a file leaves the monitored host, but its source IP
may be a proxy, VPN, NAT gateway, or security service.

\section{Methodology}

The implementation followed an incremental engineering process:

\begin{enumerate}[label=\arabic*.]
  \item establish a minimal DOCX generation and callback flow;
  \item configure a project-local monitored directory and inherited SACL;
  \item validate Windows event 4663 and dedicated Wazuh rules;
  \item generalize registration and correlation to every active HoneyDoc;
  \item add XLSX and structured formats with explicit token policies;
  \item integrate the official Canarytokens provider;
  \item preserve raw and normalized evidence in SQLite;
  \item add a unified dashboard and automated regression tests;
  \item perform manual Office, Canary email, and end-to-end Wazuh tests.
\end{enumerate}

\section{Contributions}

The main contributions of JANUS are:

\begin{itemize}
  \item a multi-format generation pipeline that separates narrative generation
  from artifact assembly;
  \item a token policy that reflects real activation semantics instead of
  promising an automatic callback for every file;
  \item an Office beacon implementation compatible with an external OOXML image
  relationship and requiring no macro;
  \item a local detection path from SACL and event 4663 to Wazuh and a JANUS
  verdict;
  \item a chain of evidence linking raw alert hashes, normalized signals, and
  HoneyDoc registry entries;
  \item a dashboard that visibly separates Wazuh observations from remote
  Canarytokens callbacks.
\end{itemize}

\section{Report organization}

Chapter 2 reviews the relevant deception, auditing, Wazuh, Canarytokens, and AI
concepts. Chapter 3 defines requirements and presents the architecture. Chapter
4 details the implementation and validation. Chapter 5 analyzes evidence
quality, limitations, security hardening, and future work.

\section{Conclusion}

The project is motivated by a simple principle: a decoy is useful only when its
interaction is observable and its evidence is interpreted honestly. The next
chapter positions JANUS among the technologies that implement this principle.

